\documentclass[11pt,a4paper]{report}

\usepackage[T1]{fontenc}
\usepackage[utf8]{inputenc}
\usepackage[magyar]{babel}
\usepackage{lmodern}
\usepackage{geometry}
\usepackage{setspace}
\usepackage{graphicx}
\usepackage{float}
\usepackage{booktabs}
\usepackage{longtable}
\usepackage{array}
\usepackage{xcolor}
\usepackage{listings}
\usepackage{hyperref}
\usepackage{enumitem}
\usepackage{caption}
\newcolumntype{L}[1]{>{\raggedright\arraybackslash}p{#1}}
\renewcommand{\arraystretch}{1.08}

\geometry{left=2.4cm,right=2.2cm,top=2.0cm,bottom=2.0cm}
\setstretch{1.15}
\setlength{\parindent}{0.8cm}
\setlength{\parskip}{0.08em}
\setlength{\textfloatsep}{10pt plus 2pt minus 2pt}
\setlength{\floatsep}{8pt plus 2pt minus 2pt}
\setlength{\intextsep}{8pt plus 2pt minus 2pt}
\setlist[itemize]{nosep,leftmargin=1.5em}
\setlist[enumerate]{nosep,leftmargin=1.7em}
\setlength{\LTpre}{4pt}
\setlength{\LTpost}{4pt}

\usepackage{titlesec}
\titleformat{\chapter}[hang]
{\normalfont\large\bfseries}
{\thechapter.}{0.6em}{}
\titlespacing*{\chapter}{0pt}{0pt}{5pt}

\titleformat{\section}
{\normalfont\normalsize\bfseries}
{\thesection}{0.55em}{}
\titlespacing*{\section}{0pt}{5pt}{2pt}

\titleformat{\subsection}
{\normalfont\normalsize\bfseries}
{\thesubsection}{0.5em}{}
\titlespacing*{\subsection}{0pt}{4pt}{2pt}


\hypersetup{colorlinks=true,linkcolor=black,urlcolor=blue,citecolor=black}

\lstdefinestyle{code}{
    basicstyle=\ttfamily\small,
    breaklines=true,
    frame=single,
    numbers=left,
    numberstyle=\tiny,
    xleftmargin=1em,
    framexleftmargin=0.5em,
    showstringspaces=false,
    tabsize=2
}
\lstset{style=code}

\newcommand{\projectname}{ProjectFlow}
\newcommand{\todo}[1]{\textcolor{red}{[TODO: #1]}}

\begin{document}

    \begin{titlepage}
        \centering
        {\Large Szegedi Tudományegyetem\par}
        \vspace{0.4cm}
        {\large Természettudományi és Informatikai Kar\par}
        \vfill
        {\Huge\bfseries SZAKDOLGOZAT\par}
        \vspace{1.2cm}
        {\LARGE\bfseries \projectname\\ Webes projekt- és ügyfélkezelő rendszer\par}
        \vfill
        \begin{minipage}{0.45\textwidth}
            \raggedright
            \textbf{Készítette:}\\
            Somorjai Lili\\
            Programtervező Informatikus BSc
        \end{minipage}
        \hfill
        \begin{minipage}{0.45\textwidth}
            \raggedleft
            \textbf{Témavezető:}\\
            Dr. Pengő Edit\\

        \end{minipage}
        \vfill
        {\large Szeged\\2026\par}
    \end{titlepage}

    \tableofcontents
    \clearpage

    \chapter{Bevezetés}

    \section{Motiváció}

    A \projectname{} ötlete a korábbi webfejlesztői tapasztalataimból indult ki.
    Az elmúlt időszakban olyan projekteken dolgoztam, ahol vállalkozások számára
    készítettünk weboldalakat, és a fejlesztés során rendszeresen több különböző
    szereplőnek kellett együttműködnie. Egy tipikus projektben jelen volt az ügyfél,
    a projektvezető, egy vagy több fejlesztő, valamint esetenként külső
    alvállalkozók is.

    A gyakorlatban azt tapasztaltam, hogy a projekt információi gyakran több,
    egymástól független rendszerben jelentek meg. A feladatok egy projektkezelő
    felületen szerepeltek, a részletes egyeztetés e-mailben vagy chatalkalmazásban
    történt, a fájlok külön tárhelyen voltak, az ügyfél visszajelzései pedig sokszor
    nem közvetlenül ahhoz a feladathoz kapcsolódtak, amelyről szóltak.

    Ez adta az alapötletet egy olyan rendszer létrehozásához, amely ugyanazon
    alkalmazáson belül képes kezelni az ügyfeleket, projekteket, feladatokat,
    projektcsapattagokat, mérföldköveket, kommenteket és jóváhagyási folyamatokat.

    \section{A dolgozat célja}

    A dolgozat célja a \projectname{} megtervezési és fejlesztési folyamatának
    bemutatása. A hangsúly nem csupán a kész funkciók felsorolásán van, hanem azon
    is, hogyan változott a rendszer a fejlesztés során, milyen problémák merültek
    fel, és ezek milyen architekturális vagy adatmodellbeli döntésekhez vezettek.

    A fő célok:
    \begin{itemize}
        \item több szervezet adatainak elkülönített kezelése;
        \item szervezeti és projektszintű szerepkörök támogatása;
        \item ügyfelek és ügyfél-kapcsolattartók kezelése;
        \item projektek, mérföldkövek és feladatok létrehozása;
        \item ügyfél és belső csapat eltérő jogosultságainak kezelése;
        \item belső és ügyfél számára látható kommentek támogatása;
        \item feladatok jóváhagyási folyamatának megvalósítása;
        \item a fontos műveletek aktivitásnaplóban történő rögzítése;
        \item a frontend és a backend üzleti logikájának szétválasztása.
    \end{itemize}

    \chapter{A rendszer fő funkciói}

    \section{Szervezetek és szerepkörök}

    A rendszer központi egysége a szervezet (\texttt{Organization}). Egy
    felhasználó egy szervezethez az \texttt{OrganizationMember} kapcsolaton
    keresztül tartozhat.

    A szervezeti szerepkörök:
    \begin{itemize}
        \item \texttt{OWNER},
        \item \texttt{ADMIN},
        \item \texttt{PROJECT\_MANAGER},
        \item \texttt{MEMBER},
        \item \texttt{CONTRACTOR},
        \item \texttt{CLIENT}.
    \end{itemize}

    A projekt szintjén külön szerepkör-rendszert alkalmaztam. Erre azért volt
    szükség, mert ugyanaz a felhasználó szervezeti és projektszinten eltérő
    felelősséggel rendelkezhet.

    \section{Ügyfélkezelés}

    A \texttt{Client} modell egy ügyfélcéget reprezentál. Ettől külön választottam
    a \texttt{ClientContact} modellt, amely egy konkrét kapcsolattartót ír le.
    A kapcsolat szándékosan opcionális a \texttt{User} felé, így egy kapcsolattartó
    akkor is eltárolható, ha még nincs saját \projectname{} fiókja.

    \section{Projektkezelés}

    Minden projekt egy szervezethez és egy ügyfélhez kapcsolódik. A projekthez
    tartozik projektvezető, projektcsapat, ügyfél-kapcsolattartók, mérföldkövek,
    feladatok és aktivitásnapló.

    \section{Feladatkezelés}

    A feladatok Kanban nézetben jelennek meg. A státuszok:
    \begin{itemize}
        \item \texttt{TODO},
        \item \texttt{IN\_PROGRESS},
        \item \texttt{REVIEW},
        \item \texttt{AWAITING\_APPROVAL},
        \item \texttt{BLOCKED},
        \item \texttt{DONE}.
    \end{itemize}

    A feladatokhoz prioritás, határidő, felelős, mérföldkő és
    \texttt{clientVisible} jelző tartozhat.

    \chapter{Felhasznált technológiák}

    \section{Next.js és React}

    A frontend és a backend is Next.js alapokra épül, de két külön alkalmazásban.
    Az \texttt{apps/web} a felhasználói felületet, az \texttt{apps/api} pedig a
    HTTP API-t valósítja meg. A két alkalmazás külön porton fut, így fejlesztés
    közben a böngészőből érkező kérések tényleges HTTP kommunikáción keresztül
    érik el a backendet.

    A frontend az App Router felépítését használja. A szerverkomponensek olyan
    adatokat kérnek le, amelyek az oldal első megjelenítéséhez szükségesek, míg a
    klienskomponensek kezelik az interaktív műveleteket, például a taskok
    szerkesztését, a Kanban drag-and-drop működését, a kommentelést és a
    jóváhagyási döntéseket.

    A frontend és az API szétválasztása azért volt fontos tervezési döntés, mert
    így az üzleti szabályok nem a megjelenítési réteghez kötődnek. A frontend
    feladata elsősorban az adatok megjelenítése és a felhasználói műveletek
    elküldése, míg a jogosultsági és állapotkezelési szabályok a backendben
    maradnak.

    \section{TypeScript}

    A teljes alkalmazás TypeScript alapú. A típusosság különösen hasznos a nagyobb
    adatstruktúrák, jogosultsági objektumok és a frontend és backend között
    megosztott szerződések kezelésében.

    A közös típusok egy része a \texttt{packages/contracts} workspace-ben
    található. Ide kerültek például az autentikációs request/response típusok és a
    közösen használt szerepkör-definíciók. Ez csökkenti annak esélyét, hogy a
    frontend és a backend ugyanazt az adatstruktúrát eltérően értelmezze.

    A TypeScript a jogosultsági logikánál is fontos: a taskhoz visszaadott
    permission objektum konkrét mezőkkel írja le, hogy az adott user milyen
    műveleteket hajthat végre. Így egy új jogosultság hozzáadásakor fordítási időben
    is láthatóvá válik, hol kell azt kezelni.

    \section{Tailwind CSS}

    A felhasználói felület stílusozását Tailwind CSS segítségével valósítottam meg.

    \section{Prisma ORM és PostgreSQL}

    Az üzleti adatokat PostgreSQL adatbázisban tárolom. Az alkalmazás a Prisma ORM
    segítségével kommunikál az adatbázissal.

    A Prisma séma deklaratívan írja le a modelleket, enumokat, relációkat,
    egyediségi feltételeket és indexeket. A fejlesztés során minden jelentősebb
    adatmodell-változást külön migrációban vezettem be. Ez különösen hasznos volt a
    fázisos fejlesztésnél, mert a korábbi adatok megtartása mellett lehetett új
    táblákat, enumokat és kapcsolatokat hozzáadni.

    A service réteg nem közvetlen SQL utasításokat használ. A perzisztencia a
    repository-kon és a Prisma kliensen keresztül történik. Ezzel az adatbázis
    elérése elkülönül az üzleti szabályoktól, és a lekérdezésekben explicit
    \texttt{select} és \texttt{include} beállításokkal szabályozható, hogy pontosan
    milyen adatok kerüljenek vissza.

    \section{Supabase Auth}

    A hitelesítést Supabase Auth végzi, míg az üzleti adatokhoz saját Prisma
    \texttt{User} rekord tartozik. A kettőt a \texttt{User.authUserId} mező
    kapcsolja össze.

    A backend egy provider interfészen keresztül használja a Supabase szolgáltatást,
    így az auth-specifikus kód nem terjed szét az alkalmazásban. A
    \texttt{SupabaseAuthProvider} végzi a bejelentkezést, regisztrációt, auth
    callback feldolgozását, valamint a meghívott fiókok előkészítését.

    Sikeres hitelesítés után az alkalmazás saját session tokent készít. A token a
    felhasználó provider azonosítóját, e-mail címét, nevét és lejárati idejét
    tartalmazza, majd HMAC-SHA256 aláírást kap. A token HttpOnly cookie-ban kerül
    eltárolásra, ezért kliensoldali JavaScriptből közvetlenül nem olvasható.

    \chapter{A rendszer architektúrája}

    \section{Monorepo felépítés}

    \begin{lstlisting}
projectflow/
|-- apps/
|   |-- web/
|   `-- api/
|-- packages/
|   `-- contracts/
|-- scripts/
|-- package.json
`-- package-lock.json
    \end{lstlisting}

    \section{Rétegek szétválasztása}

    A backend fejlesztése során fokozatosan az alábbi felépítés alakult ki:

    \begin{center}
        \texttt{Web $\rightarrow$ HTTP API $\rightarrow$ Service $\rightarrow$ Repository
            $\rightarrow$ Prisma $\rightarrow$ PostgreSQL}
    \end{center}

    A route fájlok szándékosan vékonyak. Feladatuk elsősorban a Next.js route és a
    megfelelő service összekapcsolása. A request feldolgozása, validációja,
    jogosultságellenőrzése és az üzleti döntések a
    \texttt{apps/api/lib/services/routes} könyvtárban találhatók.

    A repository réteg a Prisma kliens elérését központosítja. Ennek köszönhetően
    a service-eknek nem kell mindenhol közvetlenül ismerniük a Prisma
    implementációját. A jogosultsági helper-ek és a task access logika e fölött a
    réteg fölött működnek, így egy üzleti művelet előtt külön ellenőrizhető a
    felhasználó szervezeti, projekt- és erőforrásszintű hozzáférése.

    A projekt részletes oldalához például az \texttt{/overview} service egy
    összesített válaszobjektumot állít elő. Ebben a backend már előre kiszámítja a
    felhasználó jogosultságait és a task permission objektumokat, így a frontendnek
    nem kell a teljes üzleti szabályrendszert újraimplementálnia.

    \section{Hitelesítési architektúra}

    \begin{center}
        \texttt{Web $\rightarrow$ ProjectFlow API $\rightarrow$ AuthService
            $\rightarrow$ AuthProvider $\rightarrow$ Supabase Auth}
    \end{center}

    A frontend nem tartalmaz Supabase-specifikus hitelesítési logikát. A login és
    register űrlap a ProjectFlow API saját \texttt{/api/auth/*} végpontjait hívja.
    Az \texttt{AuthService} a provider interfészen keresztül kommunikál a külső
    auth szolgáltatóval, majd az identity alapján feloldja vagy létrehozza a saját
    Prisma \texttt{User} rekordot.

    Sikeres hitelesítés után a backend saját aláírt session tokent készít, amely
    HttpOnly cookie-ban kerül tárolásra. Minden védett kérésnél a backend olvassa
    és ellenőrzi ezt a tokent, majd a provider azonosító alapján betölti az aktuális
    üzleti felhasználót és annak szervezeti tagságait.

    Ez a megoldás szétválasztja a hitelesítést és az üzleti jogosultságot: a külső
    szolgáltató azt igazolja, hogy ki jelentkezett be, a ProjectFlow pedig a saját
    adatmodellje alapján dönti el, hogy az illető mit láthat vagy módosíthat.

    \chapter{A fejlesztés folyamata}

    A \projectname{} fejlesztése nem egyetlen nagy lépésben történt. A projektet
    egymásra épülő refaktorálási fázisokra bontottam. Ennek az volt az előnye, hogy
    minden nagyobb szerkezeti módosítás után külön ellenőrizni tudtam, hogy a korábbi
    funkciók továbbra is működnek-e, miközben fokozatosan lehetett bevezetni az új
    adatmodelleket és jogosultsági szabályokat.

    A fejlesztési folyamatban az egyes fázisok nem kizárólag új felhasználói
    funkciókat jelentettek. Több esetben egy korábban működő megoldást alakítottam
    át azért, mert a rendszer növekedésével már nem bizonyult elég pontosnak vagy
    biztonságosnak. A fejezet ezért időrendben mutatja be, hogyan jutott el a
    kezdeti ügyfél--projekt modelltől a jelenlegi, szervezeti és projektszintű
    jogosultságokat, kommunikációt és jóváhagyást is kezelő felépítésig.

    \section{A fejlesztési fázisok áttekintése}

    \begin{longtable}{p{0.12\textwidth}p{0.78\textwidth}}
        \toprule
        \textbf{Fázis} & \textbf{Fő eredmény} \\
        \midrule
        \endhead

        1. &
        Szervezeti alapok kialakítása, az \texttt{Organization} és
        \texttt{OrganizationMember} modellek bevezetése, a korábbi adatok
        szervezethez rendelésének előkészítése. \\

        2. &
        Projektszintű tagság és az első központi jogosultsági helper réteg
        bevezetése a \texttt{ProjectMember} és \texttt{ProjectRole} segítségével. \\

        3. &
        Valódi autentikáció bevezetése Supabase Auth segítségével, valamint a
        jogosultsági ellenőrzések bekötése a védett oldalak és API-k elé. \\

        4. &
        Teljes projektcsapat-kezelés: projekttag hozzáadása, szerepkör módosítása,
        eltávolítása és az utolsó projektvezető védelme. \\

        5. &
        Szervezeti meghívási folyamat, szerepkörök kezelése és tokenes
        meghívólinkek. \\

        6. &
        A projektmenedzsment adatmodell kibővítése státusszal, prioritással,
        dátumokkal, előrehaladással és mérföldkövekkel. \\

        6.1. &
        A CLIENT felhasználók konkrét ügyfélcéghez rendelése
        \texttt{ClientContact} segítségével, valamint a mérföldkövek teljes CRUD
        kezelése. \\

        7. &
        A Task modell refaktorálása és a Kanban-alapú feladatkezelés
        kialakítása. \\

        8. &
        A kommunikáció, aktivitásnapló, projekt-specifikus ügyfélhozzáférés,
        jóváhagyási folyamat és további jogosultsági hardening kialakítása. \\

        \bottomrule
    \end{longtable}

    \section{1. fázis -- szervezeti alapok}

    A projekt legelső adatmodellje még lényegesen egyszerűbb volt. A kezdeti
    adatbázisban a \texttt{User}, \texttt{Client} és \texttt{Project} modellek
    jelentették a rendszer magját. A projekt közvetlenül egy ügyfélhez és egy
    felelős felhasználóhoz kapcsolódott. Ez megfelelő volt az első működő
    prototípushoz, de nem tudta kezelni azt az esetet, amikor ugyanazt az
    alkalmazást több, egymástól független vállalkozás használja.

    A következő lépés ezért a szervezeti szint bevezetése volt. Létrejött az
    \texttt{Organization} modell, valamint az \texttt{OrganizationMember}
    kapcsolótábla. Ezzel egy felhasználó már nem egyszerűen a teljes rendszerben
    rendelkezett valamilyen szereppel, hanem egy konkrét szervezethez tartozott.

    Az \texttt{OrganizationRole} enum az alábbi szerepköröket vezette be:

    \begin{itemize}
        \item \texttt{OWNER},
        \item \texttt{ADMIN},
        \item \texttt{PROJECT\_MANAGER},
        \item \texttt{MEMBER},
        \item \texttt{CONTRACTOR},
        \item \texttt{CLIENT}.
    \end{itemize}

    A meglévő \texttt{Client} és \texttt{Project} rekordok
    \texttt{organizationId} mezőt kaptak. Mivel a korábbi adatok még nem
    rendelkeztek szervezeti kapcsolattal, külön bootstrap folyamat készült, amely a
    meglévő rekordokat az első szervezethez tudta rendelni.

    Ez a fázis alapozta meg a későbbi többfelhasználós működést. A legfontosabb
    tervezési döntés az volt, hogy a jogosultságokat nem közvetlenül a
    \texttt{User} modellen tárolom, hanem a felhasználó és a szervezet közötti
    kapcsolathoz rendelem.

    \section{2. fázis -- projektszintű tagság és jogosultsági alapréteg}

    A szervezeti szerepkörök önmagukban nem bizonyultak elegendőnek. Egy
    szervezetben dolgozó munkatárs nem feltétlenül vesz részt minden projektben,
    ezért szükségessé vált a projektszintű tagság bevezetése.

    Ebben a fázisban jött létre a \texttt{ProjectRole} enum és a
    \texttt{ProjectMember} modell. A meglévő projektfelelősök automatikusan
    \texttt{PROJECT\_MANAGER} projektszerepkört kaptak, és új projekt
    létrehozásakor a kiválasztott felelős is automatikusan projektvezetői
    tagságot kapott.

    Ezzel párhuzamosan létrehoztam az első központi jogosultsági helper réteget a
    \texttt{lib/permissions.ts} fájlban. Itt jelentek meg többek között a
    \texttt{canViewProject}, \texttt{canManageProject},
    \texttt{canManageProjectMembers}, \texttt{canManageClients},
    \texttt{requireProjectAccess} és \texttt{requireProjectManagement}
    függvények.

    Fontos döntés volt, hogy ezek a helper-ek már ekkor \texttt{userId}
    paraméterrel működtek, de valódi autentikáció hiányában még nem használtam őket
    biztonsági ellenőrzésként az API-kban. Egy nem hitelesített, kézzel megadott
    felhasználóazonosító ugyanis csak látszólagos védelmet jelentett volna.

    Az előző rendszerrel való kompatibilitás miatt a \texttt{Project.ownerId}
    mező ebben a fázisban még megmaradt. Az új \texttt{ProjectMember} modell már
    párhuzamosan létezett mellette, így a rendszer fokozatosan tudott átállni az
    új projekttagsági logikára.

    \section{3. fázis -- autentikáció és valódi autorizáció}

    A harmadik fázisban került be a valódi felhasználói hitelesítés. A megoldás
    Supabase Auth integrációra épült, amelyhez elkészült a regisztráció,
    bejelentkezés, kijelentkezés és az auth callback kezelése.

    Az alkalmazás saját \texttt{User} modelljét nem cseréltem le a Supabase
    felhasználói modelljére. Ehelyett a \texttt{User.authUserId} mezővel
    kapcsoltam össze a két réteget. Ha egy már létező ProjectFlow felhasználó azonos
    e-mail címmel jelentkezett be Supabase Authon keresztül, a rendszer a meglévő
    rekordhoz kapcsolta az auth azonosítót. Ennek köszönhetően a korábban létrejött
    szervezeti és projekttagságok nem vesztek el.

    Ebben a fázisban váltak tényleges biztonsági elemmé a korábban elkészített
    permission helper-ek. A projektlista és a dashboard már a bejelentkezett
    felhasználó hozzáférése alapján szűrődött, a projekt részletes oldal
    \texttt{requireProjectAccess} ellenőrzést kapott, és az ügyfél-, projekt- és
    felhasználói API-k is jogosultsági védelem mögé kerültek.

    A fejlesztés ezen pontján vált különösen fontossá az autentikáció és az
    autorizáció fogalmának szétválasztása. A Supabase azt igazolja, hogy ki a
    felhasználó, míg a ProjectFlow saját adatmodellje dönti el, hogy az adott
    felhasználó mit tehet a rendszerben.

    \section{4. fázis -- projektcsapat-kezelés}

    A \texttt{ProjectMember} adatmodell létrejötte után a következő lépés a teljes
    projektcsapat-kezelő folyamat megvalósítása volt.

    A projektvezetői felület lehetővé tette:
    \begin{itemize}
        \item szervezeti tag hozzáadását a projekthez;
        \item projektszerepkör megadását;
        \item projektszerepkör módosítását;
        \item projekttag eltávolítását.
    \end{itemize}

    Minden írási művelet előtt szerveroldali jogosultságellenőrzés történt.
    Projektcsapatot szervezeti \texttt{OWNER}, szervezeti \texttt{ADMIN}, illetve
    az adott projekt \texttt{PROJECT\_MANAGER} tagja kezelhetett.

    Ebben a fázisban jelent meg egy fontos adatintegritási szabály is: az utolsó
    projektvezetőt nem lehetett eltávolítani vagy alacsonyabb szerepkörre
    állítani. Ha egy több projektvezetős projektből a legacy
    \texttt{Project.ownerId} által jelölt vezetőt eltávolították vagy
    lefokozták, az \texttt{ownerId} automatikusan átkerült egy másik
    projektvezetőre.

    Ez jó példa volt arra, hogy egy CRUD művelet önmagában nem elegendő: az
    üzleti szabályokat az adatok módosításával együtt kell érvényesíteni.

    \section{5. fázis -- szervezeti meghívások}

    A következő fázis célja az volt, hogy ne csak már létező felhasználókat lehessen
    egy szervezethez rendelni. Ehhez létrejött az
    \texttt{OrganizationInvitation} modell.

    Egy szervezeti \texttt{OWNER} vagy \texttt{ADMIN} e-mail cím alapján tudott
    meghívást létrehozni. A meghíváshoz egyedi token és hét napos lejárati idő
    tartozott. A link megnyitása után a meghívott felhasználó regisztrálhatott vagy
    bejelentkezhetett, majd elfogadhatta a meghívást.

    A szervezeti szerepkörök kezelése is kibővült. Bevezettem az
    \texttt{OWNER} védelmét, valamint azt a szabályt, hogy egy
    \texttt{ADMIN} nem kezelhet másik \texttt{ADMIN} szerepkörű tagot.

    Szervezetből történő eltávolításkor az adott szervezethez tartozó
    projekttagságokat is meg kellett szüntetni. Ezzel elkerülhetővé vált, hogy egy
    már nem szervezeti tag felhasználó régi projektkapcsolatokon keresztül továbbra
    is hozzáférjen adatokhoz.

    \section{6. fázis -- a projektmenedzsment modell kibővítése}

    A korábbi fázisok főleg a többfelhasználós működés és a jogosultságkezelés
    alapjait építették ki. A hatodik fázisban ismét a projektmenedzsment
    funkcionalitás került előtérbe.

    A \texttt{Project} modell új mezőket kapott:
    \begin{itemize}
        \item \texttt{status};
        \item \texttt{priority};
        \item \texttt{startDate};
        \item \texttt{dueDate};
        \item \texttt{progress};
        \item \texttt{updatedAt}.
    \end{itemize}

    A projekt státusza a \texttt{PLANNING}, \texttt{ACTIVE},
    \texttt{ON\_HOLD}, \texttt{COMPLETED} és \texttt{CANCELLED} értékek
    egyike lehetett. A prioritás külön enumot kapott.

    Ebben a fázisban készült el a \texttt{Milestone} modell is. A mérföldkőhöz
    név, leírás, státusz, határidő és projektkapcsolat tartozik. A projektvezető,
    szervezeti admin és owner létrehozhatott mérföldkövet, míg a többi projekttag
    csak olvasási jogosultságot kapott.

    A jogosultsági szabályok is pontosításra kerültek. A backend API és a navigáció
    egyaránt figyelembe vette, hogy az adott szerepkör hozhat-e létre ügyfelet vagy
    projektet. Ezzel a frontend és a backend működése közelebb került egymáshoz:
    nem csak a menüpont tűnt el, hanem közvetlen API-hívással sem lehetett
    jogosulatlan műveletet végrehajtani.

    \section{6.1. fázis -- CLIENT stabilizálás és ClientContact}

    A hatodik fázis tesztelése során két hiányosság vált láthatóvá. Az egyik az
    volt, hogy a \texttt{CLIENT} szervezeti szerepkör még nem mondta meg, melyik
    konkrét ügyfélcéghez tartozik a felhasználó. A másik, hogy a mérföldkövek csak
    létrehozhatók voltak, teljes szerkesztési és törlési folyamat még nem tartozott
    hozzájuk.

    A megoldás a \texttt{ClientContact} modell bevezetése lett. A rekord
    tartalmazza:
    \begin{itemize}
        \item a kapcsolódó ügyfélcéget;
        \item az opcionális ProjectFlow felhasználót;
        \item a kapcsolattartó nevét;
        \item e-mail címét;
        \item opcionális pozícióját.
    \end{itemize}

    Az \texttt{OrganizationInvitation} modell \texttt{clientId} mezőt kapott,
    és \texttt{CLIENT} meghívás esetén kötelezővé vált az ügyfélcég kiválasztása.

    Fontos tervezési döntés volt, hogy a \texttt{Client} továbbra is a céget, a
    \texttt{ClientContact} pedig a konkrét személyt jelenti. Így egy ügyfélcéghez
    több kapcsolattartó is tartozhat, és nem szükséges minden kapcsolattartónak
    ProjectFlow-fiókkal rendelkeznie.

    A mérföldkő-kezelés ezzel párhuzamosan teljes CRUD működést kapott:
    létrehozás, szerkesztés, státusz- és határidő-módosítás, valamint törlés.

    \section{7. fázis -- Task refaktor és Kanban}

    A hetedik fázisban a feladatkezelést alakítottam át jelentősen. A korábbi
    egyszerű Task modell enum alapú státuszt és prioritást kapott, valamint
    megjelent a felelős, a létrehozó, a mérföldkő, a határidő és az
    ügyfél-láthatóság.

    A task státuszai ebben a fázisban:
    \begin{itemize}
        \item \texttt{TODO},
        \item \texttt{IN\_PROGRESS},
        \item \texttt{REVIEW},
        \item \texttt{BLOCKED},
        \item \texttt{DONE}.
    \end{itemize}

    A jogosultsági szabályok szerepkörtől függtek. Az owner, admin és az adott
    projekt projektvezetője teljes taskkezelést kapott. A
    \texttt{MEMBER} és \texttt{CONTRACTOR} a saját magára kiosztott task
    státuszát módosíthatta, míg a \texttt{CLIENT} csak a
    \texttt{clientVisible = true} feladatokat láthatta.

    A frontend ekkor kapta meg a Kanban alapú megjelenítést. A feladatok státusz
    szerinti oszlopokban jelentek meg, és a jogosultsággal rendelkező felhasználók
    drag-and-drop módon is mozgathatták őket.

    A projekt előrehaladása automatikusan számolhatóvá vált a lezárt taskok
    arányából:

    \begin{center}
        \texttt{progress = DONE taskok száma / összes task száma * 100}
    \end{center}

    A migráció során a meglévő feladatok megmaradtak. A korábbi string státuszok
    enum értékekké alakultak, és a régi taskok létrehozójaként a projekt korábbi
    \texttt{ownerId} felhasználója került beállításra.

    \section{8. fázis -- kommunikáció, ügyfélfolyamatok és jóváhagyás}

    A nyolcadik fázisban a cél már nem egyetlen új alapmodell létrehozása volt,
    hanem a korábban kialakított projekt- és taskrendszer valós együttműködési
    folyamatokkal történő kibővítése.

    Elsőként a taskokhoz kommentrendszer került. A
    \texttt{TaskComment} modell kétféle láthatóságot kezel:
    \texttt{INTERNAL} és \texttt{CLIENT\_VISIBLE}. Ez lehetővé tette, hogy a
    belső csapat ugyanazon a taskon belül belső megjegyzéseket írjon, miközben az
    ügyfélnek szánt kommunikáció is közvetlenül a feladathoz kapcsolódik.

    A kommentekkel együtt \texttt{ActivityLog} modell is létrejött. Az
    aktivitásnapló rögzíti a fontosabb projekt- és taskműveleteket, így a projekt
    története a felületen is visszakövethetővé vált.

    A következő probléma az ügyfél-hozzáférés pontossága volt. A
    \texttt{ClientContact} kapcsolat önmagában még túl tág hozzáférést jelentett:
    egy kapcsolattartó egy ügyfélcég több projektjéhez is hozzáférhetett volna.
    Ezért létrejött a \texttt{ProjectClientContact} modell, amely már egy konkrét
    projektet köt össze egy konkrét ügyfél-kapcsolattartóval.

    A hozzáférési modell így a következő lett:

    \begin{center}
        \texttt{Project $\rightarrow$ ProjectClientContact
            $\rightarrow$ ClientContact $\rightarrow$ User}
    \end{center}

    Ezzel az ügyféloldali hozzáférés projekt-specifikussá vált.

    A nyolcadik fázis legösszetettebb új funkciója a task-jóváhagyási folyamat
    volt. A \texttt{Task} modell \texttt{requiresApproval} mezőt kapott, a
    jóváhagyók pedig külön \texttt{TaskApproval} rekordokban jelennek meg.

    Új task státusz lett az:
    \begin{center}
        \texttt{AWAITING\_APPROVAL}
    \end{center}

    A jóváhagyási döntések:
    \begin{itemize}
        \item \texttt{PENDING},
        \item \texttt{APPROVED},
        \item \texttt{REJECTED}.
    \end{itemize}

    A folyamat tipikusan:
    \begin{center}
        \texttt{REVIEW $\rightarrow$ AWAITING\_APPROVAL
            $\rightarrow$ APPROVED $\rightarrow$ DONE}
    \end{center}

    Elutasításkor:
    \begin{center}
        \texttt{AWAITING\_APPROVAL $\rightarrow$ REJECTED
            $\rightarrow$ IN\_PROGRESS}
    \end{center}

    A jóváhagyási logika három külön műveletre lett bontva:
    \begin{itemize}
        \item a jóváhagyási konfiguráció mentése;
        \item a task jóváhagyásra küldése;
        \item a kijelölt jóváhagyó döntése.
    \end{itemize}

    A funkciók növekedésével egy másik szerkezeti probléma is előtérbe került:
    egyre több helyen kellett eldönteni, hogy az aktuális user mit tehet egy adott
    taskkal. A frontend és a különböző API service-ek párhuzamos
    jogosultságkezelése hosszú távon könnyen eltérhetett volna.

    Ezért létrejött a központi task access réteg az
    \texttt{apps/api/lib/tasks/access.ts} fájlban. A backend explicit módon
    számolja többek között a következő jogosultságokat:

    \begin{lstlisting}
canView
canEdit
canDelete
canChangeStatus
canComment
canModerateComments
canConfigureApproval
canSubmitForApproval
canDecideApproval
    \end{lstlisting}

    A frontend ezek alapján jeleníti meg az egyes vezérlőket, ugyanakkor a valódi
    biztonsági ellenőrzés továbbra is szerveroldalon történik.

    A nyolcadik fázis végén további hardening lépésekre is szükség volt.
    Szerepkörváltáskor gondoskodni kellett arról, hogy ne maradjanak vissza régi
    projekttagságok vagy task-hozzárendelések, ugyanakkor egy ügyfél
    \texttt{ClientContact} CRM-rekordja se törlődjön csak azért, mert megszűnik a
    felhasználói hozzáférése. Emiatt a CRM-adat és a ProjectFlow-fiók kapcsolata
    külön életciklust kapott.

    Az API-válaszoknál adatminimalizálást is bevezettem: a teljes
    \texttt{User} rekord helyett ahol lehet, csak az
    \texttt{id}, \texttt{name} és \texttt{email} mezők kerülnek vissza a
    frontendnek.

    \section{A fázisos fejlesztés tanulságai}

    A fázisokra bontott fejlesztés egyik fő előnye az volt, hogy az adatmodell és a
    jogosultságkezelés fokozatosan tudott alkalmazkodni a valósabb használati
    esetekhez.

    A második fázisban még elegendőnek tűnt egy általános
    \texttt{ProjectMember} kapcsolat, a hatodik fázis után azonban már külön
    ügyfél-kapcsolattartói modellre volt szükség. A hetedik fázisban a
    \texttt{clientVisible} mező elegendőnek látszott az ügyféloldali tasknézethez,
    később viszont kiderült, hogy magának a projekthez való hozzáférésnek is
    projekt-specifikusnak kell lennie.

    Hasonló fejlődés figyelhető meg a jogosultságkezelésben is. A kezdeti
    permission helper-ekből először valódi szerveroldali autorizáció, majd
    erőforrás-specifikus task permission objektum alakult ki.

    Ez a folyamat jól mutatja, hogy a rendszer architektúrája nem előre,
    változatlanul rögzített tervként készült el. Az egyes megoldásokat folyamatosan
    ellenőriztem a következő fejlesztési lépések és a manuális tesztelés során,
    majd ahol szükséges volt, az adatmodellt vagy a jogosultsági szabályokat is
    refaktoráltam.

    \chapter{A frontend fejlesztése}

    \section{Fő oldalak}

    A frontend főbb oldalai:
    \begin{itemize}
        \item bejelentkezés és regisztráció;
        \item dashboard;
        \item projektlista;
        \item új projekt létrehozása;
        \item projekt részletes oldal;
        \item ügyféllista és ügyféladatlap;
        \item szervezeti beállítások;
        \item meghívás elfogadása.
    \end{itemize}

    \section{Projekt részletes oldal}

    A projektoldal a backend \texttt{/overview} végpontjáról kapja a szükséges
    adatokat, majd több külön komponensben jeleníti meg a mérföldköveket, taskokat,
    aktivitásokat és projektcsapatot.

    \section{Kanban}

    A Kanban hat oszlopot használ:
    \begin{center}
        \texttt{TODO | IN\_PROGRESS | REVIEW | AWAITING\_APPROVAL | BLOCKED | DONE}
    \end{center}

    A frontend drag-and-drop műveleteket támogat, de minden státuszváltást a backend
    is validál.

    \section{Szerepkörfüggő felület}

    A felület az API által visszaadott jogosultságok alapján jeleníti meg a
    műveleteket. Például ügyfél nem kap task szerkesztési felületet, approval
    konfiguráció csak megfelelő jogosultsággal jelenik meg, és az ügyfél nem lát
    belső kommentet.

    \chapter{Biztonság és jogosultságkezelés}

    \section{Szerveroldali autorizáció}

    A frontend jogosultsági megjelenítése csak UX funkció. A valódi védelem
    backend oldalon történik.

    \section{Session cookie}

    A saját session token HttpOnly cookie-ban kerül tárolásra. Production
    környezetben \texttt{secure} beállítást kap és \texttt{SameSite=Lax}
    politikát használ.

    \section{CORS}

    A frontend és az API külön porton fut, ezért fejlesztési környezetben CORS
    beállítás szükséges. Az API a konfigurált \texttt{WEB\_ORIGIN} origin számára
    engedélyezi a szükséges HTTP metódusokat.

    Az API proxy az \texttt{OPTIONS} preflight kérésekre külön válaszol, és
    beállítja az \texttt{Access-Control-Allow-Origin},
    \texttt{Access-Control-Allow-Methods},
    \texttt{Access-Control-Allow-Headers} és
    \texttt{Access-Control-Allow-Credentials} fejléceket. Mivel a session cookie
    credentialként kerül továbbításra, a frontend fetch hívásai
    \texttt{credentials: "include"} beállítást használnak.

    A jóváhagyási konfiguráció bevezetésekor külön hiba forrása volt, hogy a
    \texttt{PUT} metódus kezdetben nem szerepelt az engedélyezett CORS metódusok
    között. Ez jól mutatta, hogy a frontend és backend funkcionális helyessége
    mellett az infrastruktúra-szintű HTTP konfigurációt is együtt kell kezelni.

    \section{Adatminimalizálás}

    A frontendnek csak a megjelenítéshez szükséges mezőket adja vissza az API.


    \section{E-mail provider réteg}

    A szervezeti meghívók e-mail küldését provider interfész mögé helyeztem. Az
    üzleti service-ek nem közvetlenül a Nodemailer SMTP implementációt hívják,
    hanem a központi \texttt{getEmailProvider()} factory-t.

    A jelenlegi implementáció SMTP-t használ. A kapcsolatot környezeti változók
    határozzák meg, például a host, port, secure mód, felhasználónév, jelszó és
    feladó. Ha az e-mail szolgáltatás nincs konfigurálva vagy a küldés sikertelen,
    a meghívás ettől még létrejön, és a rendszer kézzel megosztható meghívólinket
    biztosít.

    Ez a réteg azért hasznos, mert később egy másik e-mail szolgáltatóra történő
    váltásnál nem kell módosítani a meghívási üzleti logikát, csak új provider
    adaptert kell készíteni.


    \chapter{Fejlesztés közben felmerült problémák}

    \section{Frontend és backend jogosultságok eltérése}

    A funkciók számának növekedésével problémát jelentett volna, ha a frontend és
    a backend külön-külön ugyanazt a jogosultsági logikát tartja fenn. Ezt a
    backend által számított permission objektum oldotta meg.

    \section{Ügyfél-hozzáférés túl széles értelmezése}

    A korábbi ügyfél-hozzáférési modell túl sok projektet tehetett láthatóvá.
    A \texttt{ProjectClientContact} modell bevezetésével a hozzáférés
    projekt-specifikussá vált.

    \section{CRM-adat és hozzáférés összekeverése}

    A CRM kapcsolattartó törlése nem azonos a felhasználói hozzáférés
    megszüntetésével. A javított megoldás ezért csak a user-kapcsolatot bontja,
    a CRM rekordot megőrzi.

    \section{Approval és státuszok}

    A jóváhagyás bevezetése után a korábbi egyszerű
    \texttt{REVIEW $\rightarrow$ DONE} folyamat már nem volt elegendő. A
    szerveroldali státuszlogikát úgy módosítottam, hogy jóváhagyásköteles task
    csak az approval folyamat eredményeként kerülhessen \texttt{DONE} állapotba.

    \chapter{Tesztelés}

    A jelenlegi projektben a fejlesztés során elsősorban manuális funkcionális
    tesztelést végeztem.

    A fontosabb ellenőrzési esetek:
    \begin{itemize}
        \item különböző szerepkörök projekt-hozzáférése;
        \item ügyfél számára nem látható taskok elrejtése;
        \item belső kommentek elrejtése ügyfél elől;
        \item task státuszváltások;
        \item approval és rejection;
        \item meghívó elfogadása;
        \item szerepkörváltás;
        \item CRM-adatok megőrzése.
    \end{itemize}

    A repository jelenlegi állapotában még nincs teljes automatizált
    unit/integration tesztcsomag. Ez a következő fejlesztési feladatok egyike.

    \chapter{Továbbfejlesztési lehetőségek}

    \begin{itemize}
        \item automatizált unit és integration tesztek;
        \item GitHub Actions alapú CI;
        \item production deployment;
        \item részletesebb értesítési rendszer;
        \item fájlkezelés és csatolmányok;
        \item approval history;
        \item riportok és dashboard statisztikák;
        \item auditálhatóbb eseménynapló;
        \item nagy frontend komponensek további bontása;
        \item SaaS jellegű továbbfejlesztés.
    \end{itemize}

    \chapter{Összefoglalás}

    A \projectname{} fejlesztése során egy egyszerű projekt- és ügyfélkezelő
    alkalmazásból fokozatosan egy több szerepkört és összetettebb üzleti
    folyamatokat kezelő rendszer alakult ki.

    A fejlesztés egyik legfontosabb tanulsága az volt, hogy egy üzleti alkalmazás
    megbízhatóságát nem csak a látható funkciók száma határozza meg, hanem az
    adatmodell következetessége, a szerveroldali jogosultságkezelés és az üzleti
    fogalmak megfelelő szétválasztása is.

    A jelenlegi rendszer támogatja a szervezetek, ügyfelek, projektek,
    mérföldkövek és feladatok kezelését, a szerepkör-alapú hozzáférést, a
    kommentelést, a jóváhagyási folyamatot, a meghívásokat és az aktivitásnaplót.

    \appendix

    \chapter{Tervezett ábrák és képernyőképek}

    \begin{enumerate}
        \item teljes rendszerarchitektúra;
        \item egyszerűsített ER diagram;
        \item hitelesítési folyamat;
        \item jogosultsági modell;
        \item meghívási folyamat;
        \item task életciklus és approval folyamat;
        \item projektlista;
        \item projekt részletes oldal;
        \item Kanban;
        \item task részletek és kommentek;
        \item jóváhagyói nézet;
        \item szervezet- és ügyfélkezelő nézet.
    \end{enumerate}

    \chapter{Fontosabb forrásfájlok}

    A megvalósítás szempontjából legfontosabb forrásfájlokat az alábbi táblázat
    foglalja össze. Az elérési útvonalakat tördelhető formában adtam meg, hogy
    keskenyebb oldalszélességnél se folyjanak össze.

    \begin{longtable}{L{0.43\textwidth}L{0.47\textwidth}}
        \toprule
        \textbf{Fájl} & \textbf{Szerepe} \\
        \midrule
        \endhead

        \footnotesize\path{apps/api/prisma/schema.prisma}
        &
        Az adatmodell, kapcsolatok, enumok és főbb adatbázis-korlátozások definíciója.
        \\[3pt]

        \footnotesize\path{apps/api/lib/auth/service.ts}
        &
        Az autentikáció üzleti rétege, amely az auth providert és a saját User
        rekordokat kapcsolja össze.
        \\[3pt]

        \footnotesize\path{apps/api/lib/permissions.ts}
        &
        A szervezeti és projektszintű hozzáférési szabályok központi helper rétege.
        \\[3pt]

        \footnotesize\path{apps/api/lib/tasks/access.ts}
        &
        A task-specifikus jogosultságok és a frontendnek visszaadott permission
        objektumok számítása.
        \\[3pt]

        \footnotesize\path{apps/api/lib/services/routes/projects/[id]/overview/service.ts}
        &
        A projekt részletes nézetéhez szükséges összesített adatok és jogosultságok
        előállítása.
        \\[3pt]

        \footnotesize\path{apps/api/lib/services/routes/projects/[id]/tasks/[taskId]/service.ts}
        &
        A task szerkesztésének, státuszváltásának és törlésének szerveroldali szabályai.
        \\[3pt]

        \footnotesize\path{apps/api/lib/services/routes/projects/[id]/tasks/[taskId]/comments/service.ts}
        &
        A kommentek listázása és létrehozása, beleértve az ügyfél-láthatósági
        szabályokat.
        \\[3pt]

        \footnotesize\path{apps/api/lib/services/routes/projects/[id]/tasks/[taskId]/approvals/submit/service.ts}
        &
        A feladat jóváhagyásra küldésének üzleti logikája.
        \\[3pt]

        \footnotesize\path{apps/api/lib/services/routes/projects/[id]/tasks/[taskId]/approvals/[approvalId]/decision/service.ts}
        &
        A jóváhagyási és elutasítási döntések kezelése.
        \\[3pt]

        \footnotesize\path{apps/api/lib/services/routes/organizations/[id]/members/service.ts}
        &
        A szervezeti szerepkörváltások, projekttagságok és CRM-adatintegritási
        szabályok kezelése.
        \\[3pt]

        \footnotesize\path{apps/web/components/task-kanban.tsx}
        &
        A Kanban, task szerkesztés, kommentelés és approval fő frontend komponense.
        \\

        \bottomrule
    \end{longtable}

\end{document}